\paragraph{Echocardiography Video Segmentation.}
    \label{related:evs}
Echocardiography video segmentation aims for the temporally coherent segmentation of key cardiac structures, such as the left ventricle, which is crucial for the accurate assessment of cardiac functions like ejection fraction. Fully-supervised methods exhibit excellent performance, but they rely on dense, frame-by-frame annotation that is both time-consuming and costly. Sparsely-supervised methods, which require only a few annotated frames, have thus become an active area of research, aiming to effectively balance segmentation accuracy with annotation overhead. 
%
Multi-frame aggregation methods~\cite{sokeh2018superframes, liu2018path} and Space-Time Memory Networks (STM)~\cite{oh2019videoobjectsegmentationusing} are widely used in video segmentation tasks to leverage temporal information. The former processes multiple frames simultaneously to improve segmentation performance in long sequences, but demands substantial GPU resources in large-scale scenarios. By storing the features and segmentation results of past frames in a memory bank and retrieving only the necessary semantic information when a new frame arrives, STM networks significantly reduce memory overhead while maintaining temporal consistency. Methods such as XMem~\cite{cheng2022xmemlongtermvideoobject} and XMem++~\cite{bekuzarov2023xmemproductionlevelvideosegmentation} introduce multi-level memory mechanisms, drawing on human cognitive processes to sustain effective segmentation in extremely long videos while controlling memory usage.
%
Existing Space-Time Memory Networks \cite{oh2019videoobjectsegmentationusing, cheng2021rethinkingspacetimenetworksimproved, cheng2022xmemlongtermvideoobject, bekuzarov2023xmemproductionlevelvideosegmentation, cheng2024puttingobjectvideoobject} typically rely on mask propagation from fully annotated reference frames, yet they lack functionality to identify phases of the cardiac cycle. In medical applications, where it may be necessary to automatically locate end-systole and end-diastole frames, this limitation can reduce segmentation accuracy at critical moments. MemSAM~\cite{deng2024memsam} incorporates richer temporal cues into the memory bank, reducing misinterpretation of ambiguous structures and mitigating error propagation, but its computational overhead remains high, hindering direct deployment in clinical settings. Consequently, achieving both low computational cost and precise cardiac cycle recognition is a critical challenge that spatiotemporal memory methods urgently need to address in medical video segmentation.